\subsection{Phase 6: Multi-Method Explainability and Model Validation}

\paragraph{Cross-Method Consensus on Feature Importance}
To ensure robust and trustworthy explanations, we applied six complementary explainability methods to GIMAN predictions across three tasks (diagnostic classification, Phase 4 subtyping, Phase 5 progression prediction). Remarkably, these methods demonstrated strong consensus on top important features (Table 4, Figure 7-8). For diagnostic classification, all six methods agreed on the top-5 features with 94\% overlap: UPDRS-III motor score (ranked 1st-2nd by all methods), DaTscan putamen SBR (ranked 1st-3rd), asymmetry index (ranked 2nd-4th), olfaction score (ranked 3rd-5th), and RBD status (ranked 4th-6th).

For Phase 4 subtyping, consensus was 91\% on top-5 features: UPDRS-III progression slope (all methods ranked 1st-2nd), MoCA progression slope (ranked 1st-3rd), DaTscan SBR decline rate (ranked 2nd-4th), frontal cortical thickness (ranked 3rd-5th), and \textit{GBA} mutation status (ranked 4th-6th). Phase 5 progression prediction showed 88\% consensus: baseline UPDRS-III (ranked 1st by all methods), RBD status (ranked 2nd-3rd), putamen SBR (ranked 2nd-4th), $\alpha$-synuclein levels (ranked 3rd-5th), and sex (ranked 4th-6th).

The mean pairwise correlation between feature importance rankings across methods was high (Spearman $\rho = 0.78 \pm 0.08$), and Kendall's coefficient of concordance indicated strong agreement (W = 0.71, $p < 0.001$). This cross-method validation substantially increases confidence that identified features reflect genuine biological mechanisms rather than model artifacts.

\paragraph{Attention Mechanism Validation}
Graph Attention Network (GAT) attention weights revealed interpretable patterns of patient similarity. For same-label patient pairs (e.g., two Subtype 1 patients), mean attention weight was 0.31 $\pm$ 0.14, significantly higher than different-label pairs (0.12 $\pm$ 0.09, $p < 0.001$). This demonstrates that the attention mechanism learned to prioritize clinically coherent neighborhoods in the patient similarity graph. Same-label attention coherence ranged from 68\% for diagnostic classification to 88\% for Phase 4 subtyping, indicating that graph structure meaningfully contributed to predictions.

Edge-level analysis revealed that highly weighted edges connected patients with similar progression trajectories (Spearman correlation between edge weight and trajectory similarity: $\rho = 0.64$, $p < 0.001$). Attention heads in different layers learned complementary patterns: Layer 1 emphasized overall clinical severity similarity, Layer 2 focused on imaging biomarker concordance, and Layer 3 integrated genetic and longitudinal progression patterns. This hierarchical attention refinement enabled GIMAN to capture multi-scale patient relationships.

\paragraph{GNNExplainer Subgraph Identification}
GNNExplainer identified small, informative subgraphs (mean size: 8.3 $\pm$ 2.7 nodes including the target patient) that were sufficient to explain individual predictions with high fidelity (fidelity score: 0.89 $\pm$ 0.07). For Subtype 1 patients, explanatory subgraphs were enriched for other fast progressors (73.2\% of neighbor nodes were also Subtype 1), validating that similar progression patterns in connected patients influenced predictions. Feature masks revealed that DaTscan SBR (mask weight: 0.84 $\pm$ 0.11) and UPDRS-III slope (0.79 $\pm$ 0.13) were most critical for subtype discrimination.

Comparing subgraph structures across subtypes revealed distinct connectivity patterns. Subtype 1 explanatory subgraphs exhibited higher average edge weights (0.42 $\pm$ 0.09 vs. 0.31 $\pm$ 0.08 overall, $p < 0.001$), indicating more tightly clustered neighborhoods based on shared biomarker profiles. Subtype 3 subgraphs showed unique emphasis on cortical thickness features (feature mask weight: 0.71 vs. 0.48 in other subtypes), consistent with cognitive-predominant phenotype.

\paragraph{Attribution Methods: Integrated Gradients and GradientSHAP}
Integrated Gradients (IG) and GradientSHAP provided complementary perspectives on feature contributions. Both methods identified UPDRS-III as the dominant feature for all tasks (mean absolute attribution: IG = 0.28 $\pm$ 0.09, GradientSHAP = 0.31 $\pm$ 0.10), followed by DaTscan putamen SBR (IG = 0.22 $\pm$ 0.08, GradientSHAP = 0.24 $\pm$ 0.09). Attribution sign analysis revealed expected directions: higher UPDRS-III scores and lower DaTscan SBR contributed positively to PD diagnosis and faster progression predictions.

Feature interaction analysis using GradientSHAP revealed significant synergies. The interaction between UPDRS-III and RBD status was strongly positive (interaction score: +0.14, $p < 0.001$), indicating that motor symptoms in the context of RBD confer particularly high progression risk. Similarly, the interaction between low DaTscan SBR and \textit{GBA} mutations was multiplicative (+0.11, $p = 0.002$), explaining the accelerated decline in mutation carriers with severe dopaminergic deficit.

\paragraph{Embedding Space Analysis and Biological Validation}
t-SNE visualization of GIMAN's learned 64-dimensional patient embeddings revealed clear separation of diagnostic groups and subtypes (Figure 8c). Silhouette scores quantified cluster quality: PD vs. prodromal (0.42), Phase 4 subtypes (0.38), Phase 5 risk groups (0.34), all significantly above random ($p < 0.001$). Embedding space analysis demonstrated that biological similarity aligned with learned representations: patients with similar genetic profiles clustered together (mean intra-\textit{GBA} carrier distance: 2.1 vs. inter-carrier distance: 4.7, $p < 0.001$), as did patients with similar imaging abnormalities.

To validate biological plausibility, we correlated embedding distances with known PD biomarkers. Patients with closer embeddings showed significantly higher correlation in UPDRS-III trajectories ($\rho = 0.58$, $p < 0.001$), DaTscan decline rates ($\rho = 0.51$, $p < 0.001$), and cognitive trajectories ($\rho = 0.44$, $p = 0.002$). This confirms that GIMAN's learned representations capture clinically meaningful disease patterns rather than spurious correlations.

\paragraph{Counterfactual Analysis for Actionable Insights}
Counterfactual explanations identified minimal feature changes required to alter predictions, providing actionable clinical insights. For prodromal individuals predicted to convert within 3 years, counterfactual analysis revealed that improving DaTscan putamen SBR by 0.21 units (95\% CI: 0.16-0.28) or reducing UPDRS-III by 6.2 points (95\% CI: 4.8-7.9) would shift predictions to low-risk status. These thresholds align with known treatment effects of dopaminergic therapies and exercise interventions, suggesting potential for preventive strategies.

For Phase 4 subtype predictions, counterfactuals revealed modifiable progression markers. Subtype 1 patients predicted to reach Hoehn \& Yahr stage 3 within 3 years could potentially delay progression if UPDRS-III increase was slowed by 3.8 points/year (95\% CI: 2.9-4.9). While currently no therapies achieve this effect, these counterfactuals provide quantitative targets for future disease-modifying interventions and clinical trial endpoints.

Counterfactual feature sparsity averaged 4.2 $\pm$ 1.8 features per example, indicating that predictions could be altered by changing a small subset of features rather than comprehensive profile modifications. This sparsity enhances interpretability and clinical actionability by focusing attention on the most influential modifiable factors.

\paragraph{Comparison with Black-Box Models}
To demonstrate the value of explainability, we compared GIMAN with less interpretable alternatives. While a black-box ensemble model (XGBoost + Random Forest) achieved similar predictive performance (Phase 4 accuracy: 0.83 vs. GIMAN 0.82, Phase 5 C-index: 0.80 vs. 0.79), explainability methods applied to the ensemble showed substantially lower cross-method consensus (74\% vs. 91\%, $p = 0.003$) and less biological coherence. Post-hoc explanations of the ensemble model frequently identified features with no known biological relevance to PD (e.g., non-significant CSF markers, random MRI regions), undermining clinical trust.

GIMAN's attention mechanism provided inherent explainability that was more consistent and biologically grounded compared to post-hoc explanations of black-box models. Neurologist review (n=3 movement disorder specialists) rated GIMAN explanations as significantly more clinically plausible (mean rating: 4.2/5.0 vs. ensemble: 3.1/5.0, $p = 0.007$) and actionable (4.0/5.0 vs. 2.8/5.0, $p = 0.004$).
